\scanpage{80}
2. Notons que \deleted{si $\underline{\mathrm O}_S$ est de car. \unsure{$p$}} l'on tire de (1.3)
\[
(2.1)\qquad\left\{\begin{aligned}
\omega_G&=\underline M\otimes_{\mathbb Z}\underline{\mathrm O}_S\\
\underline n_G&=\underline{\operatorname{Tor}}_1^{\mathbb Z}(M,\underline{\mathrm O}_S)
\end{aligned}\right.
\]
ce qui s'écrit, si $S$ est de car. $p>0$
\[
\unsure{(2.1)|(2.2)}\qquad\left\{\begin{aligned}
\omega_G&=(M\otimes_{\mathbb Z}\mathbb F_p)\otimes_{\mathbb F_p}\underline{\mathrm O}_S=M_p\otimes_{\mathbb F_p}\underline{\mathrm O}_S\\
\underline n_G&\simeq {}_pM\otimes_{\mathbb F_p}\underline{\mathrm O}_S
\end{aligned}\right.
\]
et si $M$ est annulé par $p$, i.e. $G$ annulé par $p$ cela devient
\[
(2.2')\qquad \omega_G\simeq M\otimes_{\mathbb F_p}\underline{\mathrm O}_S,\qquad \underline n_G\simeq M\otimes_{\mathbb F_p}\underline{\mathrm O}_S.
\]
On a d'autre part une classe canonique
\[
\unsure{\xi|\zeta}\in\operatorname{Ext}^{2}_{\underline{\mathrm O}_S}(M\otimes_{\mathbb Z}\underline{\mathrm O}_S,\underline{\operatorname{Tor}}_1^{\mathbb Z}(M,\underline{\mathrm O}_S))
\]
qu'on peut interpréter, en présence d'un $\Lambda$, comme provenant d'un
\[
\unsure{\xi|\zeta}_\Lambda\in\operatorname{Ext}^{2}_{\Lambda}(M\otimes_{\mathbb Z}\Lambda,\operatorname{Tor}_1^{\mathbb Z}(M,\Lambda));
\]
Supposons p.ex. que $\Lambda=\mathbb Z/n\mathbb Z$, et que \deleted{$M\unsure{\otimes}$}\inserted{$M$} soit un $\Lambda$-module loc. libre de t.f. Alors il en est de $=$ de $M\otimes_{\mathbb Z}\Lambda\simeq M$ et de $\operatorname{Tor}_1^{\mathbb Z}(M,\Lambda)\simeq{}_nM\simeq M$, et
\[
\unsure{\xi|\zeta}_\Lambda\in\mathrm H^2(S,\operatorname{End}(M)),
\]
et \deleted{si \unsure{$\xi\in$} $\underline{\mathrm O}_S$ est plat sur $\Lambda$ (p.ex. $n=p$)}
\[
\unsure{\xi|\zeta}\in\mathrm H^2(S,\operatorname{End}(E))\qquad(E=M\otimes_\Lambda\underline{\mathrm O}_S)
\]
et $\unsure{\xi|\zeta}$ est déduit de $\unsure{\xi|\zeta}_\Lambda$ par l'hom. can.
\mspage{79}
